% =====================================================================
%  1bat-ud5-6.tex  —  SESSIÓ 6: Què és l'interès? El dilema del crèdit ràpid
%  (sense preàmbul: s'inclou des de main.tex)
% =====================================================================
\horatitol{Sessió 6 — Què és l'interès? El dilema del crèdit ràpid}%
{Un anunci de crèdit ràpid: 300 € a tornar en 30 dies pagant-ne 30 d'interessos. És barat o car? La pregunta obre la matemàtica financera.}

\subsection{Idea clau}

Amb els percentatges ben afinats (sessions 1 a 5), entrem a la part més
professionalitzadora de la unitat: la \textbf{matemàtica financera}, l'eina per
detectar abusos i assessorar famílies vulnerables sobre l'endeutament. Avui partim
d'un anunci real i ben enganyós: \emph{«$\num{300}\eur$ a l'instant; torna'ls en $30$
dies pagant només $\num{30}\eur$ d'interessos».} La pregunta del dia: \textbf{és barat
o car?}

\begin{idea}
\textbf{Què és l'interès}\\
L'\textbf{interès} és el preu que es paga per fer servir uns diners durant un temps.
\begin{itemize}[leftmargin=1.4em, itemsep=2pt]
  \item \textbf{Capital} ($C_0$): els diners prestats.
  \item \textbf{Tipus d'interès} ($r$): el percentatge que es paga \emph{per període}
        (expressat com a factor: un $10\%\rightarrow r=0,10$).
  \item \textbf{Temps} ($t$): quants períodes dura.
\end{itemize}
\textbf{Dues maneres de comptar-lo:}
\begin{itemize}[leftmargin=1.4em, itemsep=1pt]
  \item \textbf{Interès simple:} els interessos es calculen \emph{sempre sobre el
        capital inicial}.
  \item \textbf{Interès compost:} els interessos s'acumulen al capital i
        \emph{generen nous interessos} (ho formalitzarem la sessió següent).
\end{itemize}
\end{idea}

\subsection{Exemples}

\begin{exemple}[el dilema del crèdit ràpid]
$\num{300}\eur$ a tornar en $30$ dies pagant $\num{30}\eur$ d'interessos. En un mes,
els $\num{30}\eur$ sobre $\num{300}\eur$ són
\[
\frac{\num{30}}{\num{300}}=0,10=10\%.
\]
Un $10\%$\dots \emph{en un sol mes}. La trampa de l'anunci és que el $10\%$ sona poc,
però \textbf{és mensual}. Si projectem la mateixa quantitat a tot un any (dotze mesos)
amb interès simple, serien $12\per\num{30}=\num{360}\eur$ d'interessos sobre
$\num{300}\eur$ de capital: un $120\%$ anual. I si els interessos generessin
interessos (compost), encara molt més. \textbf{És caríssim.}
\end{exemple}

\begin{exemple}[interès simple: estalvi senzill]
Poso $\num{1000}\eur$ en un producte que dóna un $3\%$ \textbf{simple} anual durant
$2$ anys. Cada any genera el mateix interès, sempre sobre els $\num{1000}\eur$
inicials:
\[
\text{interès anual}=0,03\per\num{1000}=\num{30}\eur,\qquad
\text{en $2$ anys}=2\per\num{30}=\num{60}\eur.
\]
Capital final: $\num{1000}+\num{60}=\num{1060}\eur$. En interès \textbf{simple}, els
interessos \emph{no} generen interessos: cada any s'afegeixen els mateixos
$\num{30}\eur$.
\end{exemple}

\subsection{Exercicis resolts}

\begin{exresolt}[interès mensual projectat a l'any]
Un crèdit ràpid presta $\num{200}\eur$ i cobra $\num{24}\eur$ d'interessos en un mes.
Quin percentatge mensual és? I quant en interessos, amb interès simple, si durés un
any?
\solucio
\textbf{Percentatge mensual:} $\dfrac{\num{24}}{\num{200}}=0,12=12\%$ al mes.

\textbf{En un any (simple):} cada mes el mateix interès, $\num{24}\eur$, durant $12$
mesos: $12\per\num{24}=\num{288}\eur$ d'interessos sobre $\num{200}\eur$ de capital.
Això és un $\dfrac{\num{288}}{\num{200}}=1,44=144\%$ anual. \textbf{Abusiu.}
\resposta{$12\%$ mensual; $\num{288}\eur$ d'interessos en un any (un $144\%$ anual).}
\end{exresolt}

\begin{exresolt}[interès simple a uns quants anys]
Una família ingressa $\num{500}\eur$ en un dipòsit al $2\%$ simple anual. Quants
diners hi haurà al cap de $3$ anys?
\solucio
\textbf{Interès cada any:} $0,02\per\num{500}=\num{10}\eur$, sempre sobre els
$\num{500}\eur$ inicials.

\textbf{En $3$ anys:} $3\per\num{10}=\num{30}\eur$ d'interessos.

\textbf{Capital final:} $\num{500}+\num{30}=\num{530}\eur$.
\resposta{$\num{530}\eur$ (interès simple: $\num{10}\eur$ cada any).}
\end{exresolt}

\subsection{Exercicis per fer}

\begin{exercicis}
\begin{exlist}
  \item Un crèdit presta $\num{400}\eur$ i cobra $\num{40}\eur$ d'interessos en un mes.
        Quin percentatge mensual és? Quant serien els interessos, amb interès simple,
        en un any?
  \item Calcula l'interès \textbf{simple} i el capital final:
    \begin{enumerate}[label=\alph*), leftmargin=1.6em, topsep=2pt]
      \item $\num{800}\eur$ al $4\%$ simple anual durant $2$ anys;
      \item $\num{1500}\eur$ al $3\%$ simple anual durant $4$ anys.
    \end{enumerate}
  \item Un anunci diu «només un $5\%$!» però és un $5\%$ \textbf{mensual}. A quant
        equival en interessos al cap d'un any (simple) sobre un préstec de
        $\num{600}\eur$?
  \item \textbf{(Comparació)} Dos productes d'estalvi per a $\num{1000}\eur$: el primer
        dóna un $5\%$ simple durant $1$ any; el segon, un $2\%$ simple durant $3$ anys.
        Quin dóna més interessos en total?
  \item \textbf{(Interpretació)} Per què un anunci de crèdit ràpid presenta l'interès
        \textbf{mensual} en comptes de l'anual? Què intenta amagar?
  \item \textbf{(Raonament)} Explica amb les teves paraules la diferència entre interès
        simple i interès compost. En quin dels dos creixerà més de pressa un deute que
        no es paga, i per què?
\end{exlist}
\end{exercicis}

% ---------------------------------------------------------------------
\fullsolucions{Sessió 6}
\begin{solucions}
\begin{sollist}
  \item Percentatge mensual: $\dfrac{\num{40}}{\num{400}}=0,10=10\%$. En un any
        (simple): $12\per\num{40}=\num{480}\eur$ d'interessos sobre $\num{400}\eur$
        (un $120\%$ anual).
  \item (a) Interès anual $0,04\per\num{800}=\num{32}\eur$; en $2$ anys
        $\num{64}\eur$; capital final $\num{864}\eur$. \quad
        (b) Interès anual $0,03\per\num{1500}=\num{45}\eur$; en $4$ anys
        $\num{180}\eur$; capital final $\num{1680}\eur$.
  \item Un $5\%$ mensual sobre $\num{600}\eur$ són $0,05\per\num{600}=\num{30}\eur$ al
        mes; en $12$ mesos (simple), $12\per\num{30}=\num{360}\eur$ d'interessos sobre
        $\num{600}\eur$ (un $60\%$ anual).
  \item Primer: $0,05\per\num{1000}=\num{50}\eur$ en $1$ any. Segon:
        $0,02\per\num{1000}=\num{20}\eur$ l'any, en $3$ anys $\num{60}\eur$. El
        \textbf{segon} dóna més interessos en total ($\num{60}\eur$ contra $\num{50}\eur$).
  \item Perquè l'interès mensual és un número \textbf{petit} ($5\%$, $10\%$) que sona
        assumible, mentre que l'anual equivalent és enorme ($60\%$, $120\%$\dots).
        L'anunci amaga el \textbf{cost real anual} del crèdit perquè sembli barat.
  \item Resposta oberta. Idea: en l'interès \textbf{simple} els interessos es calculen
        sempre sobre el capital inicial; en el \textbf{compost} s'acumulen al capital i
        generen nous interessos. Un deute que no es paga creix més de pressa amb interès
        \textbf{compost}, perquè cada període els interessos s'afegeixen al deute i el
        període següent també generen interessos (l'«espiral»).
\end{sollist}
\end{solucions}
